\documentclass[11pt]{article}
\usepackage[margin=1in]{geometry}
\usepackage{booktabs}
\usepackage{hyperref}

\title{Fine TIME: Task-Specific Foundation-Model Adaptation}
\author{}
\date{}

\begin{document}
\maketitle

\section{Question}

Fine TIME evaluates whether task-specific adaptation of Chronos-2,
Chronos-Bolt, and TS-ICL improves over the same frozen backbone on identical
official TIME test windows. Each task is fitted only on its official training
prefix. The initial protocol uses one seed, 100 full-model optimizer steps, no
covariates, and the 90-task small profile.

\section{Inherited frozen baselines}

The executable surface contains a foundation benchmark of Chronos-2,
Chronos-Bolt, TS-ICL, and Seasonal Naive; TimesFM-3 is excluded from the active
experiment set. It also contains a Chronos-2 comparison between native
multivariate targets, independent univariate targets, and other target
histories supplied as past-only covariates. Seasonal Naive is deterministic
and provides the taskwise MASE
normalizer. It also fixes the evaluation grid to cells with finite target
support, finite Seasonal Naive predictions on that support, and finite
Seasonal Naive MASE. Every learned model uses this same grid; a non-finite
forecast on expected support invalidates its task rather than changing the
aggregation support. Inference timing excludes loading, dataset construction,
metric computation, and artifact saving.

Two further executable comparisons use the same three learned backbones and
evaluation grid. The maximum-context study evaluates five values per backbone,
halving the model-specific maximum four times. Its summary aggregates
taskwise MASE divided by the matched Seasonal Naive MASE for every model,
context-size, and horizon-size cell. The normalization study compares the
unchanged input with a per-window, per-variate z-score computed after context
truncation. Constant inputs use unit scale, and every output quantile is
transformed back to the target units before evaluation.

\section{Fine-tuning comparison}

Each model-task pair stores an independently adapted checkpoint, then evaluates
the frozen and adapted states on the shared Seasonal-defined finite grid. The
report requires exact pairs within a backbone, gives every task equal weight,
and reports mean MASE, win rates, and paired PNG/PDF figures. Artifacts live
below \texttt{outputs/<surface>/task\_finetuning/}; reports and logs remain
inside that experiment hierarchy.

\section{Evidence boundary}

Current result claims are limited to project-owned manifests and aggregates
produced under the shared Seasonal evaluation grid. No Fine TIME run has yet
completed, so the project makes no adaptation claim.

\end{document}
